\color{unidarkgreen}{\subsection{Контроль силового выключателя}}
\color{black}

\begin{enumerate}[label=\arabic{section}.\arabic{subsection}.\arabic{enumi}, labelsep=4pt, leftmargin=0pt, itemindent=57pt, itemsep=0pt, parsep=0pt] % если что можно крутить itemsep

\item Общие сведения
\begin{enumerate}[label=\arabic{section}.\arabic{subsection}.\arabic{enumi}.\arabic*, labelsep=4pt, leftmargin=0em, itemindent=65pt, parsep=0pt]
    

\item
Функциональный блок контроля силового выключателя (КСВ) предназначен для обработки 
технологической информации от выключателя и его цепей управления, а также от трансформатора тока, 
позволяя определять их состояние, информировать о возможных неисправностях и формировать 
соответствующую сигнализацию. 

\item
В случае присоединения линии к системе через два выключателя для каждого из них можно 
задействовать свою логику КСВ. Далее приведено описание КСВ для выключателя В1(ЛВ), логика КСВ 
для остальных выключателей аналогична.

\begin{figure}[H]
\centering
\includegraphics[width=1\textwidth,height=1\textheight,keepaspectratio]{\fbpath/761.3371 КСВ/img/control_Q.pdf}
\captionof{figure}{Функциональная схема контроля исправности выключателя}\label{pict:control_Q}
\end{figure}

\end{enumerate}

\item
Контроль целостности цепей управления осуществляется по сигналам состояния цепей первой и 
второй группы электромагнитов отключения (ЭМО1 и ЭМО2) и цепи электромагнита включения (ЭМВ). 
Наличие второй группы электромагнитов отключения определяется положением программного 
переключателя  \textbf{«Кол-во ЭМО В1(ЛВ)»}. При одновременном наличии или отсутствии сигналов контроля цепей ЭМО и 
ЭМВ в течение регулируемой выдержки времени \textbf{«Тср Неиспр. ЦУ В1(ЛВ)»} формируется сигнал 
\textit{«Неиспр. ЦУ В1(ЛВ)»}. Выдержка времени должна быть не менее времени взвода пружины привода 
выключателя. 

\item
Контроль готовности привода выключателя осуществляется по внешнему сигналу \textit{«Пруж. не 
заведены В1(ЛВ)»}. Сигнал о незаведенном состоянии пружины \textit{«Ср. Пруж. не заведены В1(ЛВ)»} 
формируется с выдержкой времени \textbf{«Тср Пруж. не заведены В1(ЛВ)»}.  

\item
Предусмотрен контроль давления элегаза в баке выключателя по внешним сигналам \textit{«Низк. давл. 
SF6 В1(ЛВ)»} и \textit{«Авар. давл. SF6 В1(ЛВ)»} от специальных датчиков давления. Соответствующие сигналы 
о низком и аварийном давлении элегаза \textit{«Ср. Низк. давл. SF6 В1(ЛВ)»} и \textit{«Ср. Авар. давл. SF6 В1(ЛВ)»}
формируются с выдержками времени, равными 1 с.

\item
Контроль положения автомата шинок питания осуществляется по  внешнему сигналу «Автомат ШП 
В1(ЛВ) ». Сигнал об отключенном положении автомата «Ср. Автомат ШП В1(ЛВ)» формируется с 
выдержкой времени выдержкой времени, равной 1 с.

\item
Обогрев выключателя контролируется по внешнему сигналу «Неиспр. обогр. В1(ЛВ)». 
Соответствующий сигнал о неисправности обогрева «Ср. Неиспр. обогр. В1(ЛВ)» формируется с 
регулируемой выдержкой времени \textbf{«Тср Неиспр. обогр. В1(ЛВ)»}. 

\item
Сигнализация о неисправности выключателя (сигнал «Неиспр. В1») срабатывает при наличии 
одного из сигналов: «Неиспр. ЦУ В1(ЛВ)», «Ср. Пруж. не заведены В1(ЛВ)», «Ср. Низк. давл. SF6 В1(ЛВ)», «Ср. Авар. давл. SF6 В1(ЛВ)», «Ср. Автомат ШП В1(ЛВ)», «Ср. Неиспр. обогр. В1(ЛВ)», а также от 
сигнала о непереключении фаз «Ср. ЗНФ В1(ЛВ)».

\begin{figure}[H]
\centering
\includegraphics[width=1\textwidth,height=1\textheight,keepaspectratio]{\fbpath/761.3371 КСВ/img/сontrol_TT.pdf}
\captionof{figure}{Функциональная схема контроля исправности ТТ}\label{pict_ksv:control_TT}
\end{figure}

\item Контроль ТТ
\begin{enumerate}[label=\arabic{section}.\arabic{subsection}.\arabic{enumi}.\arabic*, labelsep=4pt, leftmargin=0em, itemindent=65pt, parsep=0pt]

\item
В устройстве также предусмотрен контроль состояния ТТ. Контроль давления элегаза в ТТ 
осуществляется по внешним сигналам \textit{«Низк. давл. SF6 ТТ В1(ЛВ)»} и \textit{«Авар. давл. SF6 ТТ В1(ЛВ)»}. 
Соответствующие сигналы о низком и аварийном давлении элегаза \textit{«Ср. Низк. давл. SF6 В1(ЛВ)»} и \textit{«Ср. 
Авар. давл. SF6 В1(ЛВ)»} формируются с выдержками времени, равными 1 с. Обогрев ТТ контролируется 
по внешнему сигналу «Неиспр. обогр. ТТ В1(ЛВ)». Соответствующий сигнал о неисправности обогрева \textit{«Ср. Неиспр. обогр. 
ТТ В1(ЛВ)»} формируется с регулируемой выдержкой времени \textbf{«Тср Неиспр. обогр. ТТ В1(ЛВ) »}. По 
наличию одного из указанных сигналов формируется сигнал о неисправности ТТ \textit{«Неиспр. ТТ В1(ЛВ)»}. 
При вводе программного ключа \textbf{«Откл. от Авар. давл. SF6 ТТ В1(ЛВ)»} предусматривается действие 
на отключение выключателя при аварийном снижении давления элегаза в ТТ. 

\begin{figure}[H]
\centering
\includegraphics[width=1\textwidth,height=1\textheight,keepaspectratio]{\fbpath/761.3371 КСВ/img/сontrol_oper_Q.pdf}
\captionof{figure}{Функциональная схема контроля оперативного тока цепей сигнализации}\label{pict:сontrol_oper_Q}
\end{figure}

\item
Контроль отсутствия питания цепей сигнализации выключателя и ТТ осуществляется по внешним 
сигналам \textit{«Контр. опертока ЭМВ/ЭМО1 В1(ЛВ) »}, \textit{«Контр. опертока ЭМО2 В1(ЛВ) »} и \textit{«Контр. 
опертока сигн. В1»}. При состоянии «0» одного из сигналов срабатывает сигнализация о неисправности 
цепей сигнализации выключателя и ТТ (сигнал \textit{«Неиспр. сигн. В1/ТТ В1(ЛВ) »}) с выдержкой времени 1 
с.

\end{enumerate}

\begin{figure}[H]
\centering
\includegraphics[width=1\textwidth,height=1\textheight,keepaspectratio]{\fbpath/761.3371 КСВ/img/KSV_EMU.pdf}
\captionof{figure}{Функциональная схема защиты ЭМУ от длительного протекания тока}\label{pict_ksv:KSV_EMU}
\end{figure}

\item Контроль электромагнитов
\begin{enumerate}[label=\arabic{section}.\arabic{subsection}.\arabic{enumi}.\arabic*, labelsep=4pt, leftmargin=0em, itemindent=65pt, parsep=0pt]


\item
Для электромагнитов управления предусмотрена защита от длительного протекания тока. 
Контроль протекания тока через ЭМВ, ЭМО1 и ЭМО2 осуществляется при помощи токовых реле, 
положения контактов которых должны быть заведены на соответствующие входные сигналы ФСУ 
\textit{«Работа ЭМВ В1(ЛВ)»}, \textit{«Работа ЭМО1 В1(ЛВ)»}, \textit{«Работа ЭМО2 В1(ЛВ)»}. При длительном протекании 
тока по цепям ЭМО1 и ЭМВ формируется сигнал \textit{«Обест. ЭМО1 и ЭМВ В1(ЛВ)»}, действующий на 
независимый расцепитель автомата питания цепей ЭМО1 и ЭМВ. Аналогично при длительном 
протекании тока по цепи ЭМО2 формируется сигнал \textit{«Обест. ЭМО2 В1(ЛВ)»} для действия на 
независимый расцепитель автомата питания цепи ЭМО2. Регулируемые выдержки времени 
\textbf{«Тср защиты ЭМВ В1(ЛВ)»}, \textbf{«Тср защиты ЭМО1 В1(ЛВ)»} и \textbf{«Тср защиты ЭМО2 В1(ЛВ)»}
определяют время срабатывания на обесточивание ЭМВ, ЭМО1 и ЭМО2 соответственно. Кроме того, 
наличие сигнала \textit{«Ср. Авар. давл. SF6 В1(ЛВ)»} действует на обесточивание всех электромагнитов. 

\end{enumerate}

\begin{figure}[H]
\centering
\includegraphics[width=1\textwidth,height=1\textheight,keepaspectratio]{\fbpath/761.3371 КСВ/img/Blk_operation_Q.pdf}
\captionof{figure}{Функциональная схема блокировки управления выключателем}\label{pict_ksv:Blk_operation_Q}
\end{figure}

\item Блокировка отключения
\begin{enumerate}[label=\arabic{section}.\arabic{subsection}.\arabic{enumi}.\arabic*, labelsep=4pt, leftmargin=0em, itemindent=65pt, parsep=0pt]

\item
На блокировку отключения выключателя действует сигнал «Ср. Авар. давл. SF6 В1(ЛВ) ». 

\item
На блокировку включения действуют сигналы «Ср. Авар. давл. SF6 В1(ЛВ)», «Неиспр. ЦУ В1(ЛВ)», «Ср. Пруж. не заведены В1(ЛВ)», «Ср. Автомат ШП В1(ЛВ)». 

\item
Также предусмотрена блокировка отключения от внешнего сигнала \textit{«Внеш. блок. откл. В1(ЛВ)»}, 
блокировка включения от внешнего сигнала \textit{«Внеш. блок. вкл. В1(ЛВ)»} и блокировка управления (как 
включения, так и отключения) от внешнего сигнала \textit{«Внеш. блок. упр. В1(ЛВ)»}. 

\end{enumerate}

\item Реле фиксации
\begin{enumerate}[label=\arabic{section}.\arabic{subsection}.\arabic{enumi}.\arabic*, labelsep=4pt, leftmargin=0em, itemindent=65pt, parsep=0pt]

\item
В устройстве предусмотрена программная фиксация команд или включенного положения 
выключателя в зависимости от положения программного переключателя \textbf{«Реле фиксации В1(ЛВ)»}. В 
положении \textit{«РФК»} фиксируется команда включения, в положении \textit{«РФП»} – включенное положение. 
Квитирование сигнала фиксации \textit{«РФ В1(ЛВ)»} может осуществляться от сброса сигнализации или 
подачей оперативной команды на отключение. При отключении выключателя от защит и наличии 
сигнала \textit{«РФ В1(ЛВ)»} формируется сигнал о факте аварийного отключения выключателя \textit{«Сигнал 
несоотв. В1(ЛВ)»}, необходимый для работы АПВ.

\end{enumerate}

\begin{figure}[H]
\centering
\includegraphics[width=1\textwidth,height=1\textheight,keepaspectratio]{\fbpath/761.3371 КСВ/img/relay_fix.pdf}
\captionof{figure}{Функциональная схема реле фиксации}\label{pict_ksv:relay_fix}
\end{figure}


\small
\setlength{\extrarowheight}{0.05cm} % высота ячеек в таблице
\begin{longtable}
{|>{\centering\arraybackslash}m{4.5cm}|>{\centering\arraybackslash}m{4.5cm}|>{\centering\arraybackslash}m{3cm}|>{\centering\arraybackslash}m{3.5cm}|}
\caption{Перечень уставочных параметров КСВ\hfill\vspace{-0.5\baselineskip}}\label{ksv:tbl1} \\ 
\hhline{|-|-|-|-|}
\rowcolor{gray!30}
Наименование уставки & Диапазон & Шаг &  Значение по умолчанию \\ 
\hhline{|-|-|-|-|}
\endfirsthead

\caption*{\hspace{3pt}\emph{Продолжение таблицы \ref{ksv:tbl1}\hfill\vspace{-0.5\baselineskip}}} \\ % сделано по ГОСТ 2.105 п.6.8.7
\hhline{|-|-|-|-|}
\rowcolor{gray!30}
Наименование уставки & Диапазон & Шаг &  Значение по умолчанию \\  
\hhline{|-|-|-|-|}
\endhead
\hhline{|-|-|-|-|}

\endfoot
\endlastfoot

%%%%%%%%%%%%содержимое%%%%%%%%%%%%%%

\end{longtable}
\setlength{\extrarowheight}{0.0cm} % возврат параметра высота ячеек в таблице
\normalsize

\end{enumerate} 